
在说明如何用定理~\ref{thm:7-1-51}
为模空间
$\mathcal M^{\mathrm{main}}_{k+1}
(M,L,\{J_\rho\}_\rho:\beta;\operatorname{top}(\rho))$
定义仓西结构之前，
先说明该定理同斯塔谢夫
$A_\infty$ 空间理论 \textbf{[Sta63]}
的关系。

利用
$\mathcal M^{\mathrm{main}}_{k+1}$
及其胞腔分解，
可以如下定义拓扑空间 $X$
上的 $A_\infty$ 空间结构。
为描述相应的胞腔分解，
对每个
$A\in(\underline{k}^{+})^\ell$
取
$\mathcal M^{\mathrm{main}}_{\ell+1}$
的一个副本
$\mathcal M_A^{\mathrm{main}}$。
于是有
\begin{equation}
\partial\mathcal M^{\mathrm{main}}_{k+1}
=
\bigcup_{1\le i<i+\ell\le k}
\mathcal M^{\mathrm{main}}_{(i+1,\ldots,i+\ell)}
\times
\mathcal M^{\mathrm{main}}_{(1,\ldots,i,*,i+\ell+1,\ldots,k)}.
\tag{7.1.64}\label{eq:7-1-64}
\end{equation}
该分解导出 $A_\infty$ 公式，
正如定理~\ref{thm:7-1-51}
的边界分解导出式~\eqref{eq:7-1-50}。
\translatornote{原书 (7.1.64) 的
$1\le i<i+\ell\le k$
沿用 (7.1.52.1) 的索引写法；
按字面亦未包含首个输入位置的嵌入。
这里保留源式，不另造一套端点约定；
参见疑误 A-0429。}

\begin{definition}{7.1.65}\label{def:7-1-65}
（斯塔谢夫 \textbf{[Sta63]}）
称拓扑空间 $X$ 具有
$A_\infty$ 空间结构\liftIndex{I235}，
若存在一列连续映射
\[
\mathfrak M_k\colon
\mathcal M^{\mathrm{main}}_{k+1}\times X^k
\longrightarrow X,
\qquad k=2,3,\ldots,
\]
满足下述条件。
将 $\mathfrak M_k$
限制到分解~\eqref{eq:7-1-64}
的分量
\[
\bigl(
\mathcal M^{\mathrm{main}}_{(i+1,\ldots,i+\ell)}
\times
\mathcal M^{\mathrm{main}}_{(1,\ldots,i,*,i+\ell+1,\ldots,k)}
\bigr)\times X^k
\]
时，
它等于映射
\[
\begin{aligned}
&((\mathfrak x,\mathfrak y),
(x_1,\ldots,x_k))\\
&\quad\longmapsto
\mathfrak M_{k-\ell+1}
\bigl(
\mathfrak y,
(x_1,\ldots,x_i,
\mathfrak M_\ell(
\mathfrak x,(x_{i+1},\ldots,x_{i+\ell})),
x_{i+\ell+1},\ldots,x_k)
\bigr).
\end{aligned}
\]
\end{definition}

现在用 $\mathcal N_{k+1}$
定义两个 $A_\infty$ 空间之间的
$A_\infty$ 映射。
设 $X,Y$ 的 $A_\infty$ 结构
分别由
\[
\mathfrak M_k\colon
\mathcal M^{\mathrm{main}}_{k+1}\times X^k
\longrightarrow X,
\qquad
\mathfrak M'_k\colon
\mathcal M^{\mathrm{main}}_{k+1}\times Y^k
\longrightarrow Y
\]
给出。

\begin{definition}{7.1.66}\label{def:7-1-66}
从 $X$ 到 $Y$ 的
$A_\infty$ 映射\liftIndex{I236}
是一列连续映射
\[
\mathfrak F_k\colon
\mathcal N_{k+1}\times X^k
\longrightarrow Y,
\qquad k=1,2,\ldots,
\]
满足以下性质：

\noindent\textbf{(7.1.67.1)}
依定理~\ref{thm:7-1-51}
的识别，
将 $\mathfrak F_k$ 限制到
\[
\bigl(
\mathcal M^{\mathrm{main}}_{(i+1,\ldots,i+\ell)}
\times
\mathcal N_{(1,\ldots,i,*,i+\ell+1,\ldots,k)}
\bigr)\times X^k
\]
时，
它等于
\[
\begin{aligned}
&((\mathfrak x,\mathfrak y),
(x_1,\ldots,x_k))\\
&\quad\longmapsto
\mathfrak F_{k-\ell+1}
\bigl(
\mathfrak y,
(x_1,\ldots,x_i,
\mathfrak M_\ell(
\mathfrak x,(x_{i+1},\ldots,x_{i+\ell})),
x_{i+\ell+1},\ldots,x_k)
\bigr).
\end{aligned}
\]

\noindent\textbf{(7.1.67.2)}
依定理~\ref{thm:7-1-51}
的识别\translatornote{沿用定理 7.1.51 的原式：
若按 $\ell_0=1$ 逐字计算，
第一因子的标签长度为 $\ell_1-1$，
而下式首个映射却为
$\mathfrak F_{\ell_1}$，
即其输入数为 $\ell_1$；
不在此暗改索引，见疑误 A-0429。}，
将 $\mathfrak F_k$ 限制到
\[
\biggl(
\prod_{i=1}^{m}
\mathcal N_{(\ell_{i-1}+1,\ldots,\ell_i)}
\times
\bigl(\mathcal M^{\mathrm{main}}_{m+1}\bigr)'
\biggr)\times X^k
\]
时，
它等于
\[
\begin{aligned}
&\bigl(
((\mathfrak x_1,\ldots,\mathfrak x_m),
\mathfrak y),
(x_1,\ldots,x_k)
\bigr)\\
&\quad\longmapsto
\mathfrak M'_m
\bigl(
\mathfrak y,
(\mathfrak F_{\ell_1}
(\mathfrak x_1,(x_1,\ldots,x_{\ell_1})),
\ldots,
\mathfrak F_{\ell_m-\ell_{m-1}}
(\mathfrak x_m,
(x_{\ell_{m-1}+1},\ldots,x_{\ell_m})))
\bigr).
\end{aligned}
\]
\end{definition}

\begin{remark}{7.1.68}\label{rem:7-1-68}
斯塔谢夫曾告知作者，
具备定理~\ref{thm:7-1-51}
所列性质的空间，
已在他的学位论文中构造。
该论文的已发表版本 \textbf{[Sta63]}
给出了 $A_\infty$ 映射概念的简化形式
（其中考虑 $Y$ 为幺半群的情形；
亦见 \textbf{[Sug60]}，
并比较 \textbf{[BoVo73]} 第 11 页）。
作者确认，
斯塔谢夫的定义同
定义~\ref{def:7-1-66} 一致；
他也早已知道
定理~\ref{thm:7-1-51}，
惟本书给出的
$\mathcal N_{k+1}$
的几何描述，
我们认为是新的。

operad 语言有助于在包括
$A_\infty$ 结构的更一般背景下，
相容地组织所谓“一切皆同伦”
结构中出现的各种概念。
就与本书有关的数学发展而言，
把 operad 的思想
与模空间的虚基本链、
拓扑场论等结合，
也是自然的。
作者希望将来研究这样的推广
（见 \textbf{[Fuk07I]}）。
不过本书集中讨论
$A_\infty$ 结构，
并详细给出研究拉格朗日相交
弗洛尔理论所需的结果。
我们认为，
在目前阶段，
严格奠定直接相关的理论，
比在基础尚未牢固时
继续拓展更为迫切。
本书发展并确立的，
正是滤过 $A_\infty$ 代数的代数理论
及其在拉格朗日相交
弗洛尔理论中的几何实现。
事实上，
就“构造的一般思想”而言，
自 20 世纪 90 年代初以来，
这一主题已有大量扩展论述。
\end{remark}

\begin{remark}{7.1.69}\label{rem:7-1-69}
\textup{(1)}
上述定理~\ref{thm:7-1-51}
的证明在
$\mathcal N_{k+1}$
上定义了带角光滑流形结构。
但是，
式~\eqref{eq:7-1-52-1}、
\eqref{eq:7-1-52-2}
所给的分解，
并不恰好对应于
带角流形结构诱导的
$\mathcal N_{k+1}$ 边界分解。

为说明这一点，
考虑点
$\mathbf x
=(\Sigma,(z_0,\ldots,z_k),\rho)$
何时属于
$\mathcal N_{k+1}$
的余维 $d$ 角层。
记 $\rho^{-1}(0)$
为 $\Sigma$ 中时间分配为 $0$
的分支之并；
同样定义
$\rho^{-1}(1)$。
令
\[
d=
\pi_0\bigl(\rho^{-1}(0)\bigr)
+\pi_0\bigl(\rho^{-1}(1)\bigr),
\]
这里 $\pi_0$ 表示连通分支数。
我们断言 $\mathbf x$
在余维 $d$ 的角层中。
比较
式~\eqref{eq:7-1-57}
与~\eqref{eq:7-1-58}
即可证明。

\textup{(2)}
另须指出，
在命题~\ref{prop:7-1-61}
的证明中\translatornote{原书在此误引
“引理 7.1.63 的证明”；
该引理的证明在实体页 441 明言从略，
而坐标函数 $s_i$ 实际定义于
命题 7.1.61 的证明，见疑误 A-0430。}，
我们定义了坐标函数 $s_i$。
当对应奇点被平滑化时，
$s_i$ 取正值；
否则，它是含该奇点的
两个分支之时间分配的差。
因此，$s_i=0$ 处，
$\Sigma$ 的组合型改变。
不过，
将 $\mathcal N_{k+1}$
视为带角流形时，
并不把 $s_i=0$
的点视作两个层的分界。
这就是式~(7.1.52)
不同于带角流形分层的原因：
前者是后者的细分。
\end{remark}

现在回到定理~\ref{thm:7-1-51}
在模空间
\[
\mathcal M^{\mathrm{main}}_{k+1}
\bigl(
M,L,\{J_\rho\}_\rho:\beta;\operatorname{top}(\rho)
\bigr)
\]
上构造仓西结构的问题。
首先，
可以加入 $m$ 个内部标记点，
定义
\[
\mathcal M^{\mathrm{main}}_{k+1,m}
\bigl(
M,L,\{J_\rho\}_\rho:\beta;\operatorname{top}(\rho)
\bigr).
\]
仿照定义~\ref{def:7-1-42}
赋予它拓扑，
再如定理~\ref{thm:7-1-43}
证明其为豪斯多夫紧空间。

还考虑模空间
$\mathcal N_{k+1,m}$。
当目标空间 $M$ 为单点集时，
它正是上述
$\mathcal M^{\mathrm{main}}_{k+1,m}$
的特例。
利用定理~\ref{thm:7-1-44}
与定理~\ref{thm:7-1-51}
的证明，
可知
$\mathcal N_{k+1,m}$
是光滑流形。
现在定义遗忘映射
\[
\operatorname{Forget}\colon
\mathcal M^{\mathrm{main}}_{k+1,m}
\bigl(
M,L,\{J_\rho\}_\rho:\beta;\operatorname{top}(\rho)
\bigr)
\longrightarrow
\mathcal N_{k+1,m}.
\]

取模空间中的元素
\[
\bigl(
(\Sigma,\vec z,\vec z^{+}),
(u_\alpha),(\rho_\alpha)
\bigr)
\]
（见定义~\ref{def:4-6-6}）。
这里
$\vec z$ 与 $\vec z^+$
分别是边界和内部标记点集。
忘去映射 $u_\alpha$，
得到
$((\Sigma,\vec z,\vec z^+),(\rho_\alpha))$；
其中一些分支
$\Sigma_\alpha$
可能变得不稳定。
仿照定义~\ref{def:7-1-42}
之前构造
\[
\operatorname{Forget}_{\rho,\rho'}\colon
\mathcal M^{\mathrm{main}}_{k+1,m}(L;\beta)
\longrightarrow
\mathcal M^{\mathrm{main}}_{k'+1,m'}(L;\beta)
\]
的方式，
收缩这些分支，
得到
$(\Sigma',\vec z',\vec z^{+\prime})$。
接着定义其上的时间分配。
每个新分支
$\Sigma'_{\alpha'}$
都由原曲面中的若干分支
$\Sigma_{\alpha_i}$
连接而成；
令
\[
\rho'(\alpha')
=
\sup\{
\rho(\alpha_i)
\mid
\Sigma_{\alpha_i}
\text{ 属于 }\Sigma'_{\alpha'}
\}.
\]
对应于同一
$\Sigma'_{\alpha'}$
的所有
$\Sigma_{\alpha_i}$
之并是连通的，
而且其中有唯一的极大分支
（即最接近 $\Sigma_0$ 的分支）。
故 $\rho'(\alpha')$
恰是这个极大分支的时间分配。

现在为
$\mathcal M^{\mathrm{main}}_{k+1,m}$
的每个点构造仓西邻域。
设
\[
\bigl(
(\Sigma,\vec z,\vec z^+),
(u_\alpha),(\rho_\alpha)
\bigr)
\in
\mathcal M^{\mathrm{main}}_{k+1,m}
\bigl(
M,L,\{J_\rho\}_\rho:\beta;\operatorname{top}(\rho)
\bigr).
\]
如果其遗忘映像
由与
$(\Sigma,\vec z,\vec z^+)$
同胚的
$(\Sigma',\vec z',\vec z^{+\prime})$
表示，
则利用
定理~\ref{thm:7-1-51}，
第~7.1.2、
7.1.3 小节的论证
可以直接推广，
构造该点的仓西邻域；
这里不再重复。

否则，
加入一些内部标记点
$\vec z_*^+$，
使
$(\Sigma,\vec z,\widetilde{\vec z}^{+})$
稳定，
其中
$\widetilde{\vec z}^{+}
=\vec z^+\cup\vec z_*^+$。
借助这些额外标记点，
仿照第~7.1.4 小节末尾
命题~\ref{prop:7-1-1}、
\ref{prop:7-1-2}
的证明最后一步，
构造该点的仓西邻域。
所需的胶合论证
也与那里的证明相同。
至此便在
$\mathcal M^{\mathrm{main}}_{k+1,m}$
上构造了仓西结构。

\begin{remark}{7.1.70}\label{rem:7-1-70}
如注~\ref{rem:7-1-69}\textup{(1)}
所述，
$\mathcal N_{k+1}$
是带角流形。
这一事实用于上述论证，
以证明
$\mathcal M^{\mathrm{main}}_{k+1,m}$
具有带角的仓西结构。
然而，
域黎曼曲面的组合型
可以在
$\mathcal N_{k+1}$
的内部点改变，
即注~\ref{rem:7-1-69}\textup{(2)}
所述 $s_i=0$ 的点。
因此构造仓西结构时，
还须在这样的点进行胶合，
并证明此处坐标变换可微。
这可以仿照第 A1.4 节进行。

另一方面，
对本书目的而言，
模空间实际上不必具有
带角仓西结构。
要定义虚基本链，
只须模空间由具有仓西结构的
若干空间构成，
并沿边界胶合。
例如余维至少为 $2$
的层的局部行为并不重要；
换言之，
只需要类似同调流形
或伪流形的“仓西结构对应物”。

在各模空间上归纳构造多值截面时，
例如在
$\mathcal M^{\mathrm{main}}_{k+1}
(M',L';\beta;\operatorname{top}(\rho))$
上，
必须采用式~(7.1.52)
的细分，
而不是带角流形结构
诱导的分解。
因此得到的是
分段光滑多值截面，
而非整体光滑的多值截面；
这一点将在注~7.2.122 中讨论。

其他有关模空间的仓西结构
也可类似构造。
例如，
第~4.6.2 小节
式~\eqref{eq:4-6-27}
出现的
\[
\mathcal M^{\mathrm{main}}_{k+1}
\bigl(
M',L',\{J_{\rho,s}\}_{\rho,s};
\beta;\operatorname{top}(\rho),
\operatorname{twp}(s)
\bigr)
\]
是上述时间分配模空间的
单参数版本，
可用相同方法
在其上构造仓西结构。
同一小节
式~\eqref{eq:4-6-34}
出现的
\[
\mathcal M^{\mathrm{main}}_{k+1}
\bigl(
M',L',\{J_{1,s}\}_s;
\beta;\operatorname{twp}(s)
\bigr)
\]
则是
$\mathcal M^{\mathrm{main}}_{k+1}(\beta)$
的单参数版本。
与一对拉格朗日子流形
$(L^{(0)},L^{(1)})$
关联的滤过 $A_\infty$ 双模，
其几何构造所用模空间
也可同样处理。
\end{remark}
